\documentclass[11pt,reqno]{amsart}
\usepackage[tt=false]{libertine}
\usepackage{amssymb}
\usepackage{mathtools}
\usepackage[varbb]{newpxmath}
\usepackage[foot]{amsaddr}


\let\savedbigtimes\bigtimes
\let\bigtimes\relax
\usepackage{mathabx}
\let\bigtimes\savedbigtimes

\usepackage[margin=1in, bottom=1in]{geometry}
\usepackage{graphicx}
\usepackage{caption}
\usepackage{subcaption}
\usepackage{enumitem}
\usepackage{comment}
\usepackage{bbm}

\usepackage[usenames,dvipsnames]{xcolor}
\usepackage[colorlinks=true,
  linkcolor=teal!60!black,
  citecolor=PineGreen,
  urlcolor=RedViolet]{hyperref}
\hypersetup{bookmarksopen=true}

\usepackage[yyyymmdd,hhmmss]{datetime}

\usepackage{tikz}
\usetikzlibrary{decorations.pathreplacing,calligraphy}
\usetikzlibrary{calc}

\linespread{1.1}

\raggedbottom

% define theorem environments, commands, etc here
% theorem environments
\newtheorem{thm}{Theorem}[section]
\newtheorem{lem}[thm]{Lemma}
\newtheorem{ppn}[thm]{Proposition}
\newtheorem{prop}[thm]{Proposition}
\newtheorem{cor}[thm]{Corollary}
\newtheorem{fac}[thm]{Fact}
\newtheorem{conj}[thm]{Conjecture}

\theoremstyle{definition}
\newtheorem{dfn}[thm]{Definition}
\newtheorem{ex}[thm]{Example}
\newtheorem{asn}[thm]{Assumption}

\theoremstyle{remark}
\newtheorem{rmk}[thm]{Remark}

% equation environments
\numberwithin{equation}{section}
\def\beq#1\eeq{%
    \begin{equation}%
    #1%
    \end{equation}%
}
\def\baln#1\ealn{%
    \begin{align}%
    #1%
    \end{align}%
}
\def\balnn#1\ealnn{%
    \begin{align*}%
    #1%
    \end{align*}%
}

% macros
\def\lt{\left}
\def\rt{\right}
\def\fr{\frac}
\def\la{\langle}
\def\ra{\rangle}
\def\eps{\varepsilon}

\def\bA{{\boldsymbol{A}}}
\def\bB{{\boldsymbol{B}}}
\def\bC{{\boldsymbol{C}}}
\def\bD{{\boldsymbol{D}}}
\def\bE{{\boldsymbol{E}}}
\def\bF{{\boldsymbol{F}}}
\def\bG{{\boldsymbol{G}}}
\def\bH{{\boldsymbol{H}}}
\def\bI{{\boldsymbol{I}}}
\def\bJ{{\boldsymbol{J}}}
\def\bK{{\boldsymbol{K}}}
\def\bL{{\boldsymbol{L}}}
\def\bM{{\boldsymbol{M}}}
\def\bN{{\boldsymbol{N}}}
\def\bO{{\boldsymbol{O}}}
\def\bP{{\boldsymbol{P}}}
\def\bQ{{\boldsymbol{Q}}}
\def\bR{{\boldsymbol{R}}}
\def\bS{{\boldsymbol{S}}}
\def\bT{{\boldsymbol{T}}}
\def\bU{{\boldsymbol{U}}}
\def\bV{{\boldsymbol{V}}}
\def\bW{{\boldsymbol{W}}}
\def\bX{{\boldsymbol{X}}}
\def\bY{{\boldsymbol{Y}}}
\def\bZ{{\boldsymbol{Z}}}
\def\bDelta{{\boldsymbol{\Delta}}}
\def\bLambda{{\boldsymbol{\Lambda}}}
\def\bxi{{\boldsymbol{\xi}}}
\def\bomega{{\boldsymbol{\omega}}}
\def\btheta{{\boldsymbol{\theta}}}

\def\ba{{\boldsymbol{a}}}
\def\bb{{\boldsymbol{b}}}
\def\bc{{\boldsymbol{c}}}
\def\bd{{\boldsymbol{d}}}
\def\be{{\boldsymbol{e}}}
\def\bf{{\boldsymbol{f}}}
\def\bg{{\boldsymbol{g}}}
\def\bh{{\boldsymbol{h}}}
\def\bi{{\boldsymbol{i}}}
\def\bj{{\boldsymbol{j}}}
\def\bk{{\boldsymbol{k}}}
\def\bl{{\boldsymbol{\ell}}}
\def\bm{{\boldsymbol{m}}}
\def\bn{{\boldsymbol{n}}}
\def\bo{{\boldsymbol{o}}}
\def\bp{{\boldsymbol{p}}}
\def\bq{{\boldsymbol{q}}}
\def\br{{\boldsymbol{r}}}
\def\bs{{\boldsymbol{s}}}
\def\bt{{\boldsymbol{t}}}
\def\bu{{\boldsymbol{u}}}
\def\bv{{\boldsymbol{v}}}
\def\bw{{\boldsymbol{w}}}
\def\bx{{\boldsymbol{x}}}
\def\by{{\boldsymbol{y}}}
\def\bz{{\boldsymbol{z}}}
\def\bone{{\boldsymbol{1}}}

\def\bbA{{\mathbb{A}}}
\def\bbB{{\mathbb{B}}}
\def\bbC{{\mathbb{C}}}
\def\bbD{{\mathbb{D}}}
\def\bbE{{\mathbb{E}}}
\def\bbF{{\mathbb{F}}}
\def\bbG{{\mathbb{G}}}
\def\bbH{{\mathbb{H}}}
\def\bbI{{\mathbb{I}}}
\def\bbJ{{\mathbb{J}}}
\def\bbK{{\mathbb{K}}}
\def\bbL{{\mathbb{L}}}
\def\bbM{{\mathbb{M}}}
\def\bbN{{\mathbb{N}}}
\def\bbO{{\mathbb{O}}}
\def\bbP{{\mathbb{P}}}
\def\bbQ{{\mathbb{Q}}}
\def\bbR{{\mathbb{R}}}
\def\bbS{{\mathbb{S}}}
\def\bbT{{\mathbb{T}}}
\def\bbU{{\mathbb{U}}}
\def\bbV{{\mathbb{V}}}
\def\bbW{{\mathbb{W}}}
\def\bbX{{\mathbb{X}}}
\def\bbY{{\mathbb{Y}}}
\def\bbZ{{\mathbb{Z}}}
\def\bbone{{\mathbb{1}}}

\def\cA{{\mathcal{A}}}
\def\cB{{\mathcal{B}}}
\def\cC{{\mathcal{C}}}
\def\cD{{\mathcal{D}}}
\def\cE{{\mathcal{E}}}
\def\cF{{\mathcal{F}}}
\def\cG{{\mathcal{G}}}
\def\cH{{\mathcal{H}}}
\def\cI{{\mathcal{I}}}
\def\cJ{{\mathcal{J}}}
\def\cK{{\mathcal{K}}}
\def\cL{{\mathcal{L}}}
\def\cM{{\mathcal{M}}}
\def\cN{{\mathcal{N}}}
\def\cO{{\mathcal{O}}}
\def\cP{{\mathcal{P}}}
\def\cQ{{\mathcal{Q}}}
\def\cR{{\mathcal{R}}}
\def\cS{{\mathcal{S}}}
\def\cT{{\mathcal{T}}}
\def\cU{{\mathcal{U}}}
\def\cV{{\mathcal{V}}}
\def\cW{{\mathcal{W}}}
\def\cX{{\mathcal{X}}}
\def\cY{{\mathcal{Y}}}
\def\cZ{{\mathcal{Z}}}

\def\sA{{\mathscr{A}}}
\def\sB{{\mathscr{B}}}
\def\sC{{\mathscr{C}}}
\def\sD{{\mathscr{D}}}
\def\sE{{\mathscr{E}}}
\def\sF{{\mathscr{F}}}
\def\sG{{\mathscr{G}}}
\def\sH{{\mathscr{H}}}
\def\sI{{\mathscr{I}}}
\def\sJ{{\mathscr{J}}}
\def\sK{{\mathscr{K}}}
\def\sL{{\mathscr{L}}}
\def\sM{{\mathscr{M}}}
\def\sN{{\mathscr{N}}}
\def\sO{{\mathscr{O}}}
\def\sP{{\mathscr{P}}}
\def\sQ{{\mathscr{Q}}}
\def\sR{{\mathscr{R}}}
\def\sS{{\mathscr{S}}}
\def\sT{{\mathscr{T}}}
\def\sU{{\mathscr{U}}}
\def\sV{{\mathscr{V}}}
\def\sW{{\mathscr{W}}}
\def\sX{{\mathscr{X}}}
\def\sY{{\mathscr{Y}}}
\def\sZ{{\mathscr{Z}}}

\DeclareMathOperator{\EE}{{\bbE}}
\DeclareMathOperator{\PP}{{\bbP}}


% for comments
\newcommand{\TODO}[1]{\noindent\textup{\textsf{\color{red}[Todo: #1]}}}

% macros
\def\de{{\mathsf{d}}}
\def\disc{{\mathsf{disc}}}
\def\Ent{{\mathsf{Ent}}}

\title{One standard deviation does not suffice: A lower bound for discrepancy}

\author{}

\address[]{}
\email[]{}

\date{\today}

\begin{document}

\begin{abstract}
    In this short note, we prove that for each sufficiently large $n \in \mathbb N$ admitting an $n\times n$ Hadamard matrix, there exists $A=A_n \in \{ \pm 1 \}^{n \times n}$ such that
    \[
        \disc(A) \ge (1+\delta)\sqrt n, \quad \delta=10^{-20}, \quad \disc(A):= \inf_{ x \in \{ \pm 1 \}^n } \| Ax \|_{\infty} \,.
    \]
    In particular, this implies that
    \[
        \limsup_{n \to \infty} \sup_{ A \in \{ \pm 1 \}^{n \times n} } \frac{1}{\sqrt n} \disc(A) \ge 1+\delta >1 \,,
    \]
    thereby proving \cite[Conjecture~12]{bandeira2026randomstrasse101}. The crux of our construction is a probabilistic argument based on a random modification of Hadamard matrix. The main result of this paper were generated using GPT-6 Astra.
\end{abstract}

\maketitle

\setcounter{tocdepth}{1}
\tableofcontents

\section{Introduction and main results}
\label{s:intro}

Given $n \in \bbN$ and $A \in \bbR^{n \times n}$, we define the discrepancy of $A$ as $\disc(A):=\inf_{x\in\{\pm1\}^n}\|Ax\|_{\infty}$.
One of the motivations of this question is to understand how much smaller is $\disc(A)$ compared to the value, e.g. in expectation, when $x$ is drawn uniformly from the hypercube. In 1985, Joel Spencer \cite{spencer1985six} showed the celebrated ``Six Deviations Suffice’’ Theorem, which states that for any positive integer $n$ and any $A \in \{ \pm 1 \}^{n \times n}$ we have $\disc(A) \le 6\sqrt n$. There have since been improvements on this constant, but the question we will pursue here concerns lower bounds. In fact, as discussed in \cite{bandeira2026randomstrasse101}, there seems to be no known construction (or proof of existence) of an infinite family achieving a lower bound strictly larger than one. This motivated them formulate the following conjecture (see Conjecture~12 therein).
\begin{conj}{\label{conj:main}}
    We have
    \[
        \limsup_{n \to \infty} \sup_{ A \in \{ \pm 1 \}^{n \times n} } \frac{1}{\sqrt n} \disc(A) \ge 1+\delta >1 \,,
    \]
\end{conj}
The main goal of this short note is to prove Conjecture~\ref{conj:main}. Recall that a matrix $H_n \in \{ \pm 1 \}^{n \times n}$ is a Hadamard matrix if $H_n^{\top} H_n =n I_n$.
\begin{thm}{\label{thm:main}}
    There exists an absolute constant $N_0$ such that the following holds. For each even $n \ge N_0$ admitting an $n\times n$ Hadamard matrix $H_n$, there exists $A_n \in \{ \pm 1 \}^{n \times n}$ such that $\disc(A_n) \ge (1+\delta)\sqrt n$, where $\delta=10^{-20}>0$.
\end{thm}
In particular, since it is well known that an $n \times n$ Hadamard matrix exists for all $n=2^{k}, k \ge 1$, Theorem~\ref{thm:main} immediately implies Conjecture~\ref{conj:main}.

Our proof is based on a probabilistic argument that randomly modifies the signs of the Hadamard matrix $H_n$. More precisely, given an $n \times n$ Hadamard matrix $H$, independently for each row $i$, we choose a uniform subset $S_i\subseteq[n]$ of cardinality $r=\lfloor\delta n\rfloor$. Set
\begin{equation}{\label{eq:def-matrix-A}}
    A_{i,j}=
    \begin{cases}
        H_{i,j}, & j \notin S_i \,; \\
        -H_{i,j}, & j \in S_i \,.
    \end{cases}
\end{equation}
It suffices to show the following result.
\begin{thm}{\label{thm:prob-bound}}
    For all sufficiently large $n$ and every $n\times n$ Hadamard matrix $H$, the random matrix $A$ defined in \eqref{eq:def-matrix-A} with $\delta=10^{-20}$ satisfies
    \[
        \bbP\bigl(\disc(A)\le(1+\delta)\sqrt n\bigr) \le e^{-n/1600}+e^{-(1-\log 2)n} \,.
    \]
\end{thm}
It is clear that Theorem~\ref{thm:prob-bound} implies Theorem~\ref{thm:main}.

\subsection*{Statement on AI use}

The main results in this paper were generated by simply prompting GPT-6 Astra with Conjecture~\ref{conj:main}, without any mathematical input. 


\section{Proof of the main theorem}

In this section, we prove Theorem~\ref{thm:prob-bound}. We first record a number of probabilistic and combinatorial inputs that will be used in the proof. The proof of those inputs are postponed to Section~3.
We first record a bound for the number of ``nearly flat'' images of Hadamard matrices.
\begin{lem}\label{lem:count}
    Let $H\in\{-1,1\}^{n\times n}$ satisfy $H^{\top}H=nI$, where $n\ge2$. Then
    \[
        \#\mathcal F(H):=\#\left\{ x\in\{-1,1\}^n: \left\|\frac{1}{\sqrt n} Hx\right\|_1> \frac9{10}n \right\} \le 2^n e^{-n/800} \,.
    \]
\end{lem}
We then record two facts for a random $r$-subset $S \subseteq [n]$.
\begin{lem}\label{lem:tail}
    Let $1\le r\le n$ be an integer, let $w_1,\ldots,w_n\in\{-1,1\}$, let $S$ be a uniform $r$-subset of $[n]$, and let $B=\sum_{j\in S}w_j$. Then, for $u\ge0$,
    \[
         \bbP(B-\bbE B\ge u)\le e^{-u^2/(2r)}.
    \]
    The same estimate holds for the lower tail.
\end{lem}
\begin{lem}\label{lem:clt}
    Fix $\delta\in(0,1)$ and a compact interval $\Gamma\subset\bbR$.
    Let $r=\lfloor\delta n\rfloor$, and suppose that
    $w\in\{-1,1\}^n$ satisfies $\sum_jw_j=v\sqrt n$ with
    $v\in\Gamma$. Let $S$ be a uniform $r$-subset of $[n]$, and put
    $Y=n^{-1/2}(\sum_jw_j-2\sum_{j\in S}w_j)$.
    There exists $\varepsilon_n=\varepsilon_n(\Gamma,\delta)\ge0$
    tending to zero such that, uniformly over all such $w$,
    \[
        \sup_{T\in\bbR}\left|
        \bbP(Y\le T)-
        \Phi\left(\frac{T-(1-2\delta)v}
        {2\sqrt{\delta(1-\delta)}}\right)
        \right|\le\varepsilon_n.
    \]
    Here $\Phi$ is the standard normal distribution function.
\end{lem}

For $v\in\bbR$ such that
$v\sqrt n\in\{-n,-n+2,\ldots,n\}$, choose
$w\in\{-1,1\}^n$ with $\sum_jw_j=v\sqrt n$ and define
\[
    p_n(v)=\bbP\left(
        \left|\sum_jw_j-2\sum_{j\in S}w_j\right|
        \le(1+\delta)\sqrt n\right),
    \qquad |S|=\lfloor\delta n\rfloor,
\]
where $S$ is uniformly chosen. This probability depends only on $v$,
not on the particular choice of $w$.

\begin{lem}\label{lem:envelope}
    Set $\delta=10^{-20}$, $\eta=1/1600$, and
    $\lambda=1/\sqrt{8\pi\delta}$. For all sufficiently large $n$
    and every $v\in\bbR$ satisfying
    $v\sqrt n\in\{-n,-n+2,\ldots,n\}$, we have
    \begin{equation}\label{eq:envelope}
        p_n(v)\le
        \exp\left(-\log2+\eta-\lambda(v^2-1)\right).
    \end{equation}
\end{lem}

We now provide the proof of Theorem~\ref{thm:prob-bound}.
\begin{proof}[Proof of Theorem~\ref{thm:prob-bound}]
    Put $t=1+\delta$, $\delta_n=r/n$, and $a_n=1-2\delta_n$, and take $n$ to be sufficiently large that $r\ge1$ and Lemma~\ref{lem:envelope} applies. Thus $\delta_n\le\delta$ and $a_n\ge1-2\delta$. We use $\eta=1/1600$ and $\lambda=1/\sqrt{8\pi\delta}$ as in Lemma~\ref{lem:envelope}. We will apply a union bound to control $\|Ax\|_{\infty}$ for all $x\in\{\pm1\}^n$. Fix a sign vector $x$ and write $v=Hx/\sqrt n$. Orthogonality gives
    \begin{equation}\label{eq:energy}
        \sum_i v_i^2=n.
    \end{equation}
    We divide into two cases:

    \medskip
    {\em Case~1}: $x\in\mathcal F(H)$. The random choices in different rows of $A$ are independent, so Lemma~\ref{lem:envelope} and \eqref{eq:energy} imply the uniform bound
    \[
    \begin{aligned}
        \bbP(\|Ax\|_\infty\le t\sqrt n)
        &=\prod_i p_n(|v_i|)\\
        &\le\exp\left(( -\log2+\eta)n-\lambda\sum_i(v_i^2-1)\right)
        =2^{-n}e^{\eta n}.
    \end{aligned}
    \]
    Combined with Lemma~\ref{lem:count} and a union bound, this gives
    \begin{equation}\label{eq:prob-case-1}
    \begin{aligned}
        &\bbP\left(\|Ax\|_{\infty}\le t\sqrt n
            \mbox{ for some }x\in\mathcal F(H)\right)\\
        &\hspace{20mm}\le2^n e^{-n/800}\cdot2^{-n}e^{\eta n}
        =e^{-n/1600}.
    \end{aligned}
    \end{equation}

    \medskip
    {\em Case~2}: $x\notin\mathcal F(H)$. Then $\sum_i|v_i|\le0.9n$. Lemma~\ref{lem:tail} gives
    \begin{equation}\label{eq:rowtail}
    \begin{aligned}
        \bbP(\|Ax\|_\infty\le t\sqrt n)
        &=\prod_i p_n(|v_i|)\\
        &\le\prod_{i=1}^n\exp\left(-\frac{(a_n|v_i|-t)_+^2}{8\delta_n}\right).
    \end{aligned}
    \end{equation}
    Set $u_i=(a_n|v_i|-t)_+$. For $w\ge0$ one has $w^2\le tw+w(w-t)_+$. Consequently, by Cauchy--Schwarz and \eqref{eq:energy},
    \[
    \begin{aligned}
        a_n^2n
        &\le ta_n\sum_i|v_i|+a_n\sum_i|v_i|u_i\\
        &\le\frac9{10}ta_n n
             +a_n\sqrt n\left(\sum_i u_i^2\right)^{1/2}.
    \end{aligned}
    \]
    Since $a_n-9t/10\ge1/10-29\delta/10>0$, division by $a_n\sqrt n$, rearrangement, and squaring yield
    \[
        \sum_i u_i^2\ge\left(a_n-\frac9{10}t\right)^2n
        \ge\left(\frac1{10}-\frac{29}{10}\delta\right)^2n.
    \]
    Substituting this estimate into \eqref{eq:rowtail} and applying a union bound, we obtain
    \begin{equation}\label{eq:prob-case-2}
    \begin{aligned}
        &\bbP\left(\|Ax\|_{\infty}\le t\sqrt n
            \mbox{ for some }x\notin\mathcal F(H)\right)\\
        &\hspace{20mm}\le2^n\exp\left(-\frac{(1/10-(29/10)\delta)^2n}{8\delta}\right)
        \le e^{-(1-\log2)n}.
    \end{aligned}
    \end{equation}
    Combining \eqref{eq:prob-case-1} and \eqref{eq:prob-case-2} proves the result.
\end{proof}



\section{Preliminary inputs}

In this section, we complete the omitted proofs in Lemmas~\ref{lem:count}--\ref{lem:envelope}.

\subsection{Proof of Lemma~\ref{lem:count}}

\begin{proof}
Let $Q=H/\sqrt n$, let $X$ be uniform on the sign cube, and put
$F(x)=\|Qx\|_1$. Orthogonality of any two columns of $H$ implies
that $n$ is even. Write $n=2m$. An elementary binomial summation,
together with $\binom{2m}{m}/4^m\le1/\sqrt{\pi m}$, gives
\[
\begin{aligned}
    \bbE F(X)
    =\sqrt n\,\bbE\left|\sum_{j=1}^n X_j\right|
     =\sqrt n\,2m\,\frac{\binom{2m}{m}}{4^m} \le\sqrt{\frac2\pi}\,n<\frac45n.
\end{aligned}
\]
We next prove the concentration estimate needed here. Denote by $x^{(j)}$ the vector obtained by reversing the $j$th coordinate of $x$. Choose $y_i\in\{-1,1\}$ such that $y_i(Qx)_i=|(Qx)_i|$; either sign may be used when $(Qx)_i=0$. The vector $g=Q^{\top}y$ is a subgradient of $F$ at $x$, and $\|g\|_2^2=n$. Convexity therefore gives
\[
    F(x)-F(x^{(j)})\le 2x_jg_j, \qquad \sum_{j=1}^n (F(x)-F(x^{(j)}))_+^2\le4n \,.
\]
Talagrand's concentration inequality then implies that
\[
 \bbP(F(X)\ge\bbE F(X)+u)\le e^{-u^2/(8n)}.
\]
Taking $u=n/10$ proves the claimed counting bound.
\end{proof}

\subsection{Proof of Lemma~\ref{lem:tail}}

\begin{proof}
If $r=n$, then $B$ is deterministic and the conclusion is immediate.
Suppose that $r<n$. Generate $S$ by sampling distinct indices
$J_1,\ldots,J_r$ uniformly without replacement. Let
$\mathcal G_k=\sigma(J_1,\ldots,J_k)$ and
$M_k=\bbE[B\mid\mathcal G_k]$, so that $M_0=\bbE B$ and $M_r=B$.
For $1\le k\le r$, let
$m_{k-1}=(n-k+1)^{-1}\sum_{j\notin\{J_1,\ldots,J_{k-1}\}}w_j$
be the mean of the remaining signs. A direct computation gives
$M_k-M_{k-1}=(n-r)(w_{J_k}-m_{k-1})/(n-k)$.
Conditionally on $\mathcal G_{k-1}$, this increment has mean zero
and is supported on an interval of length at most
$2(n-r)/(n-k)\le2$.

For any centered random variable $Z$ supported on an interval of
length at most two, one has $\bbE e^{\theta Z}\le e^{\theta^2/2}$
for every $\theta\in\bbR$. Indeed, the second derivative of
$\log\bbE e^{\theta Z}$ is the variance under the corresponding
exponentially tilted law, which is at most one because its support
remains in the same interval. Integrating twice and using
$\bbE Z=0$ proves the bound. Applying this conditionally to the
martingale increments and iterating gives
$\bbE e^{\theta(B-\bbE B)}\le e^{r\theta^2/2}$.
For $u>0$, Markov's inequality with $\theta=u/r$ yields
\[
    \bbP(B-\bbE B\ge u)
    \le\exp\left(-\theta u+\frac{r\theta^2}{2}\right)
    =\exp\left(-\frac{u^2}{2r}\right).
\]
The case $u=0$ is immediate. Replacing every $w_j$ by $-w_j$
gives the lower-tail estimate.
\end{proof}

\subsection{Proof of Lemma~\ref{lem:clt}}

\begin{proof}
Let $K=\#\{j:w_j=1\}=(n+v\sqrt n)/2$, put
$p=K/n$ and $\delta_n=r/n$, and let
$X=\#\{j\in S:w_j=1\}$. Then $p=1/2+v/(2\sqrt n)$ and
\[
    \bbP(X=k)
    =\frac{\binom Kk\binom{n-K}{r-k}}{\binom nr},
    \qquad
    Y=(1-2\delta_n)v-\frac4{\sqrt n}(X-rp).
\]
Throughout the proof, all error estimates are uniform over the
allowed $v\in\Gamma$. Since $\Gamma$ is compact,
$p=1/2+O_\Gamma(n^{-1/2})$, while $\delta_n\to\delta$.
In particular, $p$ and $\delta_n$ stay bounded away from zero and
one for all sufficiently large $n$.

Put $\tau_n^2=n\delta_n(1-\delta_n)p(1-p)$.
We first establish a uniform local normal approximation. Fix $M>0$
and consider integers $k$ with $|k-rp|\le M\sqrt n$.
Writing $u=k-rp$, Stirling's formula gives
\[
\begin{aligned}
    \log\bbP(X=k)
    ={}&K h\left(\delta_n+\frac uK\right)
       +(n-K)h\left(\delta_n-\frac u{n-K}\right)
       -nh(\delta_n)\\
       &\quad-\frac12\log(2\pi\tau_n^2)+o(1),
\end{aligned}
\]
where $h(s)=-s\log s-(1-s)\log(1-s)$. The error is uniform
for the indicated $k$ and $v$. Taylor expansion about $\delta_n$
cancels the linear terms, and $h''(s)=-1/[s(1-s)]$ gives
\[
\begin{aligned}
    &K h\left(\delta_n+\frac uK\right)
       +(n-K)h\left(\delta_n-\frac u{n-K}\right)
       -nh(\delta_n)\\
    &\hspace{35mm}
       =-\frac{u^2}{2\tau_n^2}
        +O_{\Gamma,\delta,M}(n^{-1/2}).
\end{aligned}
\]
Consequently,
\[
    \bbP(X=k)
    =\frac{1+o(1)}{\sqrt{2\pi}\,\tau_n}
       \exp\left(-\frac{(k-rp)^2}{2\tau_n^2}\right),
    \qquad |k-rp|\le M\sqrt n,
\]
uniformly over $k$ and $v$.

The sampling indicators give $\bbE X=rp$ and
$\operatorname{Var}(X)=rp(1-p)(n-r)/(n-1)=\tau_n^2n/(n-1)$.
Thus Chebyshev's inequality bounds
$\bbP(|X-rp|>M\sqrt n)$ by $C_{\Gamma,\delta}/M^2$,
uniformly in $v$. Summing the local approximation over
$|k-rp|\le M\sqrt n$ and comparing with the corresponding
Gaussian Riemann sums, uniformly in the upper endpoint, therefore
yields
\[
    \sup_{z\in\bbR}
    \left|
       \bbP\left(\frac{X-rp}{\tau_n}\le z\right)-\Phi(z)
    \right|\longrightarrow0,
\]
uniformly over all allowed $v$.
Here one first lets $n\to\infty$ for fixed $M$ and then lets
$M\to\infty$; the Gaussian tails also vanish uniformly because
$\tau_n/\sqrt n$ stays bounded above and away from zero.

Finally, uniformly over $v\in\Gamma$,
\[
    (1-2\delta_n)v\longrightarrow(1-2\delta)v,
    \qquad
    \frac{4\tau_n}{\sqrt n}
    =4\sqrt{\delta_n(1-\delta_n)p(1-p)}
    \longrightarrow2\sqrt{\delta(1-\delta)}.
\]
Using the representation of $Y$ above, reflection symmetry of the
standard normal distribution, and continuity of its distribution
function proves the asserted uniform estimate. More explicitly,
the uniform distribution-function approximation also holds for
left limits, so reflecting $(X-rp)/\tau_n$ causes no endpoint
problem. The convergence of the means and of the positive standard
deviations then gives uniform convergence of the normal distribution
functions for all $T\in\bbR$.
\end{proof}

\subsection{Proof of Lemma~\ref{lem:envelope}}

\begin{proof}
Replacing $w$ by $-w$ shows that $p_n(-v)=p_n(v)$, so it is enough
to consider $v\ge0$. Put $t=1+\delta$, $a=1-2\delta$, and
$\sigma=2\sqrt{\delta(1-\delta)}$. Also write
$\delta_n=\lfloor\delta n\rfloor/n$ and $a_n=1-2\delta_n$,
and take $n$ sufficiently large that $\lfloor\delta n\rfloor\ge1$.
Put $B=\sum_{j\in S}w_j$ and $Y=v-2B/\sqrt n$. Then
$\bbE B=\delta_n v\sqrt n$ and
$Y=a_nv-2(B-\bbE B)/\sqrt n$.
Since $\{|Y|\le t\}\subseteq\{Y\le t\}$, Lemma~\ref{lem:tail}
gives
\[
    p_n(v)
    \le\exp\left(-\frac{(a_nv-t)_+^2}{8\delta_n}\right)
    \le\exp\left(-\frac{(av-t)_+^2}{8\delta}\right),
\]
where the second inequality uses $a_n\ge a$ and $\delta_n\le\delta$.

We also use the elementary Gaussian inequality
$\log\Phi(z)\le-\log2+\sqrt{2/\pi}\,z$ for every $z\in\bbR$.
To see this, let $\varphi(z)=(2\pi)^{-1/2}e^{-z^2/2}$ and
$R(z)=\varphi(z)/\Phi(z)$. Differentiation gives
$R'(z)=-R(z)(z+R(z))<0$, since
$z\Phi(z)+\varphi(z)=\int_{-\infty}^z(z-u)\varphi(u)\,du>0$.
Thus $\log\Phi$ is concave, and its tangent at zero is the claimed
upper bound.

\medskip
\noindent\emph{Case 1: $0\le v\le1-10\sqrt\delta$.}
In this range,
$\lambda(1-v^2)\ge(20-100\sqrt\delta)/\sqrt{8\pi}>\log2$.
The right-hand side of \eqref{eq:envelope} is therefore greater
than one, and the conclusion follows from $p_n(v)\le1$.

\medskip
\noindent\emph{Case 2: $|v-1|\le10\sqrt\delta$.}
Apply Lemma~\ref{lem:clt} with the fixed compact interval
$\Gamma=[1-10\sqrt\delta,1+10\sqrt\delta]$. Uniformly over the
allowed $v$ in this interval,
\[
    p_n(v)\le\bbP(Y\le t)
    \le\Phi\left(\frac{t-av}{\sigma}\right)+\varepsilon_n,
    \qquad \varepsilon_n\longrightarrow0.
\]
The Gaussian inequality and the choice
$\lambda=1/\sqrt{8\pi\delta}$ give
\[
\begin{aligned}
    \log\Phi\left(\frac{t-av}{\sigma}\right)
       +\log2+\lambda(v^2-1)
    &\le \sqrt{\frac2\pi}\frac{t-av}{\sigma}
       +\lambda(v^2-1)\\
    &=\lambda(v-1)^2
      +\frac{t-av-\sqrt{1-\delta}(1-v)}
      {\sqrt{2\pi\delta(1-\delta)}}.
\end{aligned}
\]
The numerator in the last fraction equals
$1+\delta-\sqrt{1-\delta}
 +v(\sqrt{1-\delta}-1+2\delta)$, which lies between zero and
$2\delta(1+v)\le6\delta$. Moreover,
$\lambda(v-1)^2\le25\sqrt\delta$ and
$\sqrt{2\pi\delta(1-\delta)}\ge2\sqrt\delta$.
The last display is consequently at most
$30\sqrt\delta<\eta/2$. Hence
\[
    p_n(v)\le
    \exp\left(-\log2+\frac\eta2-\lambda(v^2-1)\right)
    +\varepsilon_n.
\]
Since $e^{-\log2-\lambda(v^2-1)}\ge e^{-6}$ throughout this
interval, taking $n$ so large that
$\varepsilon_n\le e^{-6}(e^\eta-e^{\eta/2})$ proves
\eqref{eq:envelope} uniformly in this case.

\medskip
\noindent\emph{Case 3: $v\ge1+10\sqrt\delta$.}
Put $u=v-1\ge10\sqrt\delta$. Since $\delta=10^{-20}$,
one has $av-t=(1-2\delta)u-3\delta\ge3u/4$ and
$\lambda\le1/(4\sqrt\delta)$. Therefore
\[
\begin{aligned}
    \frac{(av-t)^2}{8\delta}-\lambda(v^2-1)
    &\ge\frac{9u^2}{128\delta}
         -\frac{u^2+2u}{4\sqrt\delta}\\
    &\ge\frac{u^2}{16\delta}-\frac{u}{2\sqrt\delta}
     \ge\frac54>\log2.
\end{aligned}
\]
The second inequality uses $\sqrt\delta\le1/32$, and the third
uses $u/\sqrt\delta\ge10$. The sampling tail bound at the start
of the proof now gives
$p_n(v)\le\exp(-\log2-\lambda(v^2-1))$, which is stronger than
\eqref{eq:envelope}. The three cases exhaust all nonnegative $v$,
and symmetry completes the proof.
\end{proof}

\bibliographystyle{alpha}
\bibliography{bib}

\end{document}


